\chapter*{Resumen}
\addcontentsline{toc}{chapter}{Resumen}

El creciente consumo energético de la infraestructura computacional ha situado la eficiencia energética del software como un criterio relevante en la ingeniería de sistemas. Sin embargo, la literatura comparativa sobre estilos arquitectónicos se ha centrado mayoritariamente en el contraste entre monolitos y microservicios distribuidos, con resultados inconsistentes según las condiciones experimentales, dejando sin evidencia empírica suficiente el efecto energético de modularizar un sistema dentro de una misma unidad de despliegue. El objetivo de esta investigación fue determinar si existen diferencias en el consumo energético del software entre una arquitectura monolítica tradicional y una arquitectura de monolito modular bajo condiciones y cargas de trabajo similares.

Se desarrolló un estudio cuantitativo de tipo experimental con diseño factorial $2 \times 2$ en bloques completos aleatorizados ($n = 12$ repeticiones por condición, para un total de 48 corridas medidas y 12 mediciones del equipo en reposo). Se construyó un banco de pruebas en memoria implementado en Java con dos versiones funcionalmente equivalentes de un mismo sistema de gestión de tienda ---monolito tradicional con esquema compartido frente a monolito modular con cinco módulos, datos privados, interfaces públicas con objetos de transferencia inmutables y eventos internos síncronos--- sometidas a una carga fija de $18\,391\,515$ operaciones por hilo en dos niveles de concurrencia: normal (un hilo) y pesada (dos hilos). La energía del procesador y la memoria DRAM se midió mediante los contadores RAPL en un equipo con procesador Intel Celeron N4020 bajo Arch Linux, y las diferencias se contrastaron con las pruebas $t$ de Welch, $U$ de Mann-Whitney y prueba pareada por bloque ($\alpha = 0.05$).

Los resultados muestran que el monolito modular consumió significativamente más energía por operación que el monolito tradicional en ambas cargas ($p < 10^{-4}$ en todas las pruebas), rechazándose la hipótesis nula $H_0$. Con carga normal, el modular registró un incremento relativo de $+4.62$\,\% en energía total por operación (3.5265 frente a $3.3707~\mu\text{J/op}$) y de $+5.08$\,\% en energía neta; con carga pesada, el incremento fue de $+6.22$\,\% (2.1426 frente a $2.0171~\mu\text{J/op}$) y $+6.24$\,\% en energía neta. La potencia media permaneció prácticamente invariable entre arquitecturas (diferencias de $+0.31$\,\% y $+0.02$\,\%), por lo que el sobrecosto energético se explica por el mayor tiempo de ejecución requerido por el monolito modular ($+4.30$\,\% y $+6.20$\,\%). Se concluye que la modularización interna conlleva un costo energético medible y moderado en términos absolutos, el cual debe ponderarse junto con criterios de mantenibilidad al tomar decisiones de diseño arquitectónico.

\vspace{0.5cm}
\noindent\textbf{Palabras clave:} arquitectura de software, monolito modular, monolito tradicional, consumo energético, eficiencia energética, RAPL.